\section{Scope and conventions}
\label{sec:scope}
The central objects are sequences and polynomial equations equipped with
arithmetic coordinates. Residue projectors select addresses in a sequence,
finite differences measure compatibility between nearby addresses, and
valuation bounds control the passage from finite precision to a $p$-adic
limit. These constructions provide a common language for the Euler packet
results and the later lifting results. Each theorem nevertheless retains its
own coefficient ring, topology and nonsingularity hypotheses.

The arithmetic framework is developed together with its completed
extensions. The stated external Nishibuchi conjecture is proved in its full
specified parameter range. The dyadic refinements and polynomial deformation theorems give stronger
conclusions under the hypotheses stated with each result. The native conjectures retain their original
identifiers, 13.2 and 13.3, when discussed as research questions, although the
section numbering of this consolidated manuscript is new.

Throughout, $\N=\{0,1,2,\ldots\}$. When $p$ is prime, $\Zp$ denotes the
$p$-adic integers and $\Qp$ their fraction field. We normalize $v_p(p)=1$.
For a rational or $p$-adic quantity, a displayed lower valuation bound is
understood with the mathematical convention $v_p(0)=+\infty$.
The corresponding formal predicates are zero-safe: the quantity is zero or
its valuation meets the stated bound. An exact finite valuation assertion
requires a nonzero quantity, and this is stated where used. Congruence in
$\Zp$ means ideal divisibility. Thus
\[
 x\equiv y\pmod{p^r}\quad\Longleftrightarrow\quad
 x-y\in p^r\Zp.
\]
For vectors, ideal membership and congruence are coordinatewise.

Write $\Delta_u f(n)=f(n+u)-f(n)$ and use $\Delta_u^k$ for its $k$-fold
iterate. Euler coefficients indexed by ordinary degree and by packet row are
distinguished explicitly. A packet displacement by $u$ rows is a degree
displacement by $Nu$ in a packet of width $N$.

The main text presents theorem statements and their arguments; accompanying
source maps identify the Lean declarations and archived files that support
them. The formal companion identifies the scope of each recorded build. The
91-check polynomial record covers the polynomial development; separate
records support the arithmetic results.
The inherited formal development is retained without adding Lean declarations. The strengthening chapters add written proofs and clearly mark their formalization boundary; no historical-priority claim is made. The analytic inverse-function results
used in Section~\ref{sec:regularity} come from Mathlib.
